\documentclass[11pt,a4paper]{article}
\usepackage{fontspec,amsmath,amssymb,amsthm,microtype}
\setmainfont{lmroman10-regular.otf}[BoldFont=lmroman10-bold.otf,ItalicFont=lmroman10-italic.otf,BoldItalicFont=lmroman10-bolditalic.otf]
\setmonofont{lmmono10-regular.otf}[Scale=MatchLowercase]
\usepackage[margin=25mm]{geometry}
\usepackage[unicode,hidelinks]{hyperref}
\setlength{\emergencystretch}{3em}
\setlength{\parskip}{0.4em}
\setlength{\parindent}{0pt}
\allowdisplaybreaks[2]
\raggedbottom
\title{\texorpdfstring{Scalar Gaussian linear maps give jointly Gaussian processes}{Scalar Gaussian linear maps give jointly Gaussian processes}}
\author{LangevinLean mathematical exposition}
\date{}
\begin{document}\maketitle
\noindent Source module: \texttt{Unadjusted\allowbreak{}Langevin/\allowbreak{}Gaussian\allowbreak{}Linear\allowbreak{}Process.\allowbreak{}lean}. The accompanying HTML pairs this mathematical exposition with the exact Lean statements and records the scope of the formalization.


If a linear map \(T:H\to L^2(P;\mathbb R)\) sends every input to a Gaussian variable, all finite collections of its evaluations are jointly Gaussian.


\section{\texorpdfstring{has\allowbreak{}Gaussian\allowbreak{}Law pi of linear\allowbreak{}Map Lp}{hasGaussianLaw pi of linearMap Lp}}


Use the finite-dimensional Gaussian criterion: every scalar linear functional of \((T(f_i))_i\) must be Gaussian. A finite linear combination equals \(T(\sum_i a_if_i)\) as an \(L^2\) class by linearity, hence agrees almost everywhere as a random variable. The assumed Gaussianity of every \(T(x)\) proves the criterion.



{\raggedright\noindent Lean declarations: \texttt{has\allowbreak{}Gaussian\allowbreak{}Law\_\allowbreak{}pi\_\allowbreak{}of\_\allowbreak{}linear\allowbreak{}Map\_\allowbreak{}Lp}.\par}

\section{\texorpdfstring{is\allowbreak{}Gaussian\allowbreak{}Process of linear\allowbreak{}Map Lp}{isGaussianProcess of linearMap Lp}}


For each finite set of times, apply the preceding joint-Gaussian theorem to the corresponding finite family of vectors \(f(t)\). This verifies the finite-dimensional-law definition of \texttt{Is\allowbreak{}Gaussian\allowbreak{}Process}. The index set may be uncountable; this statement alone asserts no sample-path continuity.



{\raggedright\noindent Lean declarations: \texttt{is\allowbreak{}Gaussian\allowbreak{}Process\_\allowbreak{}of\_\allowbreak{}linear\allowbreak{}Map\_\allowbreak{}Lp}.\par}

\end{document}
